% Part:sets-functions-relations
% Chapter: sets
% Section: reduction

\documentclass[../../../include/open-logic-section]{subfiles}

\begin{document}

\olfileid{sfr}{siz}{red}

\olsection{न्यूनीकरण}

\begin{editorial}
  या विभागात \olref[nen]{sec} मधील निष्कर्षांशी जुळणारी अगणनीयता
  न्यूनीकरणाने सिद्ध केली आहे. \olref[nen-alt]{sec} मधील निष्कर्षांशी
  जुळणारी, किंचित अधिक संक्षिप्त पर्यायी आवृत्ती \olref[red-alt]{sec}
  मध्ये दिली आहे.
\end{editorial}

कर्णीकरणाच्या युक्तिवादाने $\Pow{\PosInt}$ हा !!{nonenumerable} आहे,
हे आपण दाखवले. सर्व अनंत $0$--$1$ अनुक्रमांचा संच~$\Bin^\omega$ हा
!!{nonenumerable} आहे, याची सिद्धता आपल्याकडे आधीच होती.
$\Pow{\PosInt}$ हा !!{nonenumerable} आहे हे सिद्ध करण्याचा आणखी एक मार्ग
असा: \emph{$\Pow{\PosInt}$ हा !!{enumerable} असेल, तर $\Bin^\omega$
  हाही !!{enumerable} असतो} हे दाखवा. $\Bin^\omega$ हा !!{enumerable}
नाही, हे माहीत असल्याने $\Pow{\PosInt}$ हाही तसा असू शकत नाही. यालाच
एक समस्या दुसऱ्या समस्येकडे \emph{न्यूनीत करणे} म्हणतात---या बाबतीत
$\Bin^\omega$ चे प्रगणन करण्याची समस्या आपण $\Pow{\PosInt}$ चे प्रगणन
करण्याच्या समस्येकडे न्यूनीत करतो. नंतरच्या समस्येचे उत्तर---म्हणजे
$\Pow{\PosInt}$ चे प्रगणन---पहिल्या समस्येचे उत्तर---म्हणजे
$\Bin^\omega$ चे प्रगणन---देईल.

संच~$B$ च्या प्रगणनाची समस्या संच~$A$ च्या प्रगणनाच्या समस्येकडे कशी
न्यूनीत करायची? $A$ चे प्रगणन $B$ च्या प्रगणनात रूपांतरित करण्याची
पद्धत आपण देतो. त्यासाठीचा सर्वांत सोपा मार्ग म्हणजे
$f\colon A \to B$ असे !!a{surjective} फलन परिभाषित करणे. $x_1$, $x_2$,
\dots{} हे~$A$ चे प्रगणन असेल, तर $f(x_1)$, $f(x_2)$, \dots{} हे~$B$
चे प्रगणन करील. आपल्या बाबतीत आपण
$f\colon \Pow{\PosInt} \to \Bin^\omega$ असे आच्छादक फलन शोधत आहोत.

\begin{prob}
$g\colon B \to A$ असे !!{injective} फलन असेल आणि $B$ हा
!!{nonenumerable} असेल, तर $A$ हाही अगणनीय असतो, हे दाखवा.
$g$ वापरून~$A$ चे प्रगणन~$B$ च्या प्रगणनात कसे रूपांतरित करता येते,
हे दाखवून सिद्धता द्या.
\end{prob}

\begin{proof}[न्यूनीकरणाने {\olref[nen]{thm:nonenum-pownat}} ची सिद्धता]
$\Pow{\PosInt}$ हा !!{enumerable} आहे, आणि म्हणून त्याचे $Z_{1}$,
$Z_{2}$, $Z_{3}$, \dots{} असे प्रगणन आहे, असे समजा.

$f \colon \Pow{\PosInt} \to \Bin^\omega$ हे फलन पुढीलप्रमाणे परिभाषित
करा: $f(Z)$ हा असा अनुक्रम~$s$ असू द्या की $n \in Z$ असेल तेव्हा
आणि केवळ तेव्हाच $s(n) = 1$, आणि अन्यथा $s(n) = 0$. याने
फलन स्पष्टपणे परिभाषित होते, कारण $Z \subseteq \PosInt$ असताना कोणताही
$n \in \PosInt$ हा $Z$ चा !!a{element} असतो किंवा नसतो. उदाहरणार्थ, धन सम संख्यांचा संच
$2\PosInt = \{2, 4, 6, \dots\}$ हा $010101\dots$ या अनुक्रमाकडे,
रिक्त संच $0000\dots$ कडे आणि $\PosInt$ हा संच स्वतः $1111\dots$ कडे
पाठवला जातो.
%READERNOTE{OLSIZ-004: गोठवलेल्या इंग्रजी व्याख्येत f(Z) ला s_k असे म्हटले आहे, पण k चा Z शी संबंध किंवा निवड दिलेली नाही. अभिप्रेत अनुक्रम Z नेच ठरतो; मराठी व्याख्येत एकच बद्ध निर्गत नाव s आणि त्याचे घटक s(n) सातत्याने वापरले आहेत.}

हे फलन !!{surjective} सुद्धा आहे: $0$ आणि $1$ यांचा प्रत्येक अनुक्रम धन
पूर्णांकांच्या कोणत्या तरी संचाशी जुळतो; तो संच म्हणजे अनुक्रमात ज्या
स्थानी~$1$ आहे त्या स्थानांना अनुरूप पूर्णांकांचा संच. अधिक नेमकेपणाने,
$s \in \Bin^\omega$ असे समजा. पुढीलप्रमाणे $Z \subseteq \PosInt$
परिभाषित करा:
\[
Z = \Setabs{n \in \PosInt}{s(n) = 1}
\]
मग $f$ च्या व्याख्येवरून तपासता येते की $f(Z) = s$.

आता पुढील यादी विचारात घ्या:
\[
f(Z_1), f(Z_2), f(Z_3), \dots
\]
$f$ हे !!{surjective} असल्याने $\Bin^\omega$ मधील प्रत्येक सदस्य हा
$f$ च्या कोणत्या तरी फलसाधकावरील मूल्य असला पाहिजे, आणि म्हणून तो
यादीत आला पाहिजे. त्यामुळे ही यादी $\Bin^\omega$ मधील सर्व घटकांचे
प्रगणन करते.

म्हणून $\Pow{\PosInt}$ हा !!{enumerable} असता, तर $\Bin^\omega$ हाही
!!{enumerable} असता. पण $\Bin^\omega$ हा !!{nonenumerable} आहे
(\olref[nen]{thm:nonenum-bin-omega}). त्यामुळे $\Pow{\PosInt}$ हा
!!{nonenumerable} आहे.
\end{proof}

\begin{explain}
न्यूनीकरण कोणत्या दिशेने केले जाते याबद्दल गोंधळ होणे सोपे आहे.
उदाहरणार्थ, $g \colon \Bin^\omega \to B$ असे !!a{surjective} फलन
$B$ हा !!{nonenumerable} आहे, हे \emph{सिद्ध करत नाही}. ($g
\colon \Bin^\omega \to \Bin$ हे $g(s) = s(1)$ ने परिभाषित केलेले,
$0$ आणि $1$ यांच्या अनुक्रमाला त्याच्या पहिल्या !!{element} कडे
पाठवणारे फलन विचारात घ्या. काही अनुक्रम $0$ ने आणि काही $1$ ने सुरू
होत असल्याने ते !!{surjective} आहे. पण $\Bin$ हा सांत संच आहे.)
तसेच फलन~$f$ हे !!{surjective} असलेच पाहिजे; अन्यथा युक्तिवाद लागू
होत नाही: $f(x_1)$, $f(x_2)$, \dots{} यात~$B$ मधील सर्व !!{element}s
असतील याची खात्री उरत नाही. उदाहरणार्थ,
\[
h(n) = \underbrace{000\dots0}_{\text{$n$ वेळा $0$}}111\dots
\]
याने $h\colon \PosInt \to
\Bin^\omega$ असे फलन परिभाषित होते: प्रत्येक मूल्याची सुरुवात n शून्यांनी
होते आणि त्यानंतर कायम एक येतात. या फलनाच्या व्याप्तीत अखेरीस कायम एक
न होणारे अनुक्रम येत नाहीत, म्हणून ते आच्छादक नाही; पण $\PosInt$ हा
!!{enumerable} आहे.
%READERNOTE{OLSIZ-005: गोठवलेल्या इंग्रजी प्रदर्शनात h(n) चे मूल्य फक्त n शून्यांची सांत चिन्हमाला आहे, जी अनंत अनुक्रमांच्या Bin-omega संचात येत नाही. मराठी प्रदर्शनात त्यानंतर अनंत एकांची शेपटी जोडून वैध, अनाच्छादक फलन दिले आहे.}
\end{explain}

\begin{prob}
धन पूर्णांकांच्या जोड्यांच्या \emph{सर्व संचांचा} संच न्यूनीकरणाच्या
युक्तिवादाने !!{nonenumerable} आहे, हे दाखवा.
\end{prob}

\begin{prob}\label{sfr:siz:red:prob:nat-nat}
  $f\colon \Nat \to \Nat$ अशा सर्व फलनांचा संच~$X$ हा न्यूनीकरणाच्या
  युक्तिवादाने !!{nonenumerable} आहे, हे दाखवा. (सूचना: $X$ पासून
  $\Bin^\omega$ कडे जाणारे आच्छादक फलन द्या.)
\end{prob}

\begin{prob}
नैसर्गिक संख्यांच्या सर्व अनंत अनुक्रमांचा संच~$\Nat^\omega$ हा
न्यूनीकरणाच्या युक्तिवादाने !!{nonenumerable} आहे, हे दाखवा.
\end{prob}

\begin{prob}
$P$ हा धन पूर्णांकांच्या संचापासून $\{0\}$ या संचाकडे जाणाऱ्या फलनांचा
संच असू द्या; आणि $Q$ हा धन पूर्णांकांच्या संचापासून $\{0\}$ या
संचाकडे जाणाऱ्या \emph{अंशतः} फलनांचा संच असू द्या. $P$ हा
!!{enumerable} आहे आणि $Q$ नाही, हे दाखवा. (सूचना: $\Bin^\omega$ चे
प्रगणन करण्याची समस्या~$Q$ चे प्रगणन करण्याच्या समस्येकडे न्यूनीत करा.)
\end{prob}

\begin{prob}
$S$ हा धन पूर्णांकांच्या संचापासून \{0,1\} या संचाकडे जाणाऱ्या सर्व
!!{surjective} फलनांचा संच असू द्या; म्हणजे, $S$ मध्ये सर्व
!!{surjective}~$f\colon \PosInt \to \Bin$ आहेत. $S$ हा
!!{nonenumerable} आहे, हे दाखवा.
\end{prob}

\begin{prob}
सर्व वास्तव संख्यांचा संच~$\Real$ हा !!{nonenumerable} आहे, हे दाखवा.
\end{prob}

\end{document}
